Put
\[
h_P(x):=\sum_{a\in\mathcal A}\omega_a(x)\mu_a(x),
\qquad
z_a(x):=|\omega_a(x)|\sigma_a(x),
\]
and define the compact allocation set and its pointwise objective by
\[
\mathcal E_{\varepsilon}:=
\left\{u\in[\varepsilon,1]^K:\sum_{a\in\mathcal A}u_a=1\right\},
\qquad
F_z(u):=\sum_{a\in\mathcal A}\frac{z_a^2}{u_a}.
\tag{1}
\]
The boundary condition \(0<\varepsilon<1/K\) makes \(\mathcal E_{\varepsilon}\) nonempty.  Compactness and continuity give a minimizer; deterministic lexicographic tie breaking makes the selected minimizer \(u^*(z)\) measurable.  All divisions below are valid because \(u_a\geq\varepsilon\).  Bounded weights and the variance range give \(0\leq z_a(x)\leq W\sqrt{v_{\max}}\).

We first prove the allocation inequality needed later.  Let
\[
\Phi(z):=\min_{u\in\mathcal E_{\varepsilon}}F_z(u).
\]
The map \((z,u)\mapsto\sum_a z_a^2/u_a\) is jointly convex on \(\mathbb R_+^K\times\mathcal E_{\varepsilon}\), because each summand is the perspective of the square function.  Partial minimization over the convex set \(\mathcal E_{\varepsilon}\) therefore makes \(\Phi\) convex.  The Karush--Kuhn--Tucker equations show, whenever \(z\neq0\), that
\[
u_a^*(z)=\max\{\varepsilon,z_a/\lambda\},
\qquad
\sum_{a\in\mathcal A}u_a^*(z)=1
\tag{2}
\]
for the unique normalizing \(\lambda>0\).  Thus the optimizer is unique on every positive coordinate; any nonuniqueness concerns only zero coordinates, on which the derivative with respect to \(z_a\) is zero.  Danskin's elementary directional-derivative calculation consequently gives the single gradient
\[
\nabla_a\Phi(z)=\frac{2z_a}{u_a^*(z)}.
\tag{3}
\]
Equations (2)--(3) also show directly that the finitely many active-set pieces join with the same first derivative.  Hence active-set changes, including changes at a zero coordinate, create no first-order jump.

For \(z,\widetilde z\in\mathbb R_+^K\), put \(\delta=z-\widetilde z\) and \(\widetilde u=u^*(\widetilde z)\).  Expanding the square at the fixed allocation \(\widetilde u\), then using convexity of \(\Phi\), gives
\[
\begin{aligned}
F_z(\widetilde u)
&=\Phi(\widetilde z)+\nabla\Phi(\widetilde z)^{\mathsf T}\delta
  +\sum_{a\in\mathcal A}\frac{\delta_a^2}{\widetilde u_a},\\
\Phi(z)&\geq
\Phi(\widetilde z)+\nabla\Phi(\widetilde z)^{\mathsf T}\delta.
\end{aligned}
\]
Since \(u^*(z)\) minimizes \(F_z\), these displays prove the cross-evaluated regret inequality
\[
0\leq F_z\{u^*(\widetilde z)\}-F_z\{u^*(z)\}
\leq \sum_{a\in\mathcal A}\frac{(z_a-\widetilde z_a)^2}{u_a^*(\widetilde z)}
\leq\varepsilon^{-1}\|z-\widetilde z\|_2^2.
\tag{4}
\]
This argument controls all active-set crossings and does not require the optimizer itself to be Lipschitz.

We next specify an information order that guarantees conditional centering.  Let
\[
\mathcal G_{t-1}:=\sigma(O_1,\ldots,O_{t-1})
\]
augmented only by randomization already used before time \(t\), and let
\(\mathcal G_{t-1}^{+}:=\sigma(\mathcal G_{t-1},X_t)\).  The design constructed below ignores every coordinate of \(H_{t-1}\) other than the actually revealed past records.  Thus even if the ambient information convention places future baseline covariates in \(H_0\), the constructed kernel is measurable with respect to \((\mathcal G_{t-1},X_t)\) and does not use them.  Since the latent units are i.i.d., \(X_t\) is independent of \(\mathcal G_{t-1}\), and, conditional on \(\mathcal G_{t-1}^{+}\), the current potential outcomes have their superpopulation conditional law given \(X_t\).

Partition the assignment times, apart from a fixed initialization block, into geometric epochs
\[
I_j:=\{2^j,\ldots,(2^{j+1}-1)\wedge n\},
\qquad N_j:=2^j.
\tag{5}
\]
For an evaluation epoch \(I_j\), use the block \(I_{j-2}\) to fit the means and the subsequent disjoint block \(I_{j-1}\) to fit the variances.  Both block sizes are within fixed factors of \(N_j\).

For completeness, we record the expected-risk calculation under predictable inverse-propensity sampling.  On a cube of side \(b\), use the rescaled tensor-polynomial basis of degree \(\lfloor u\rfloor\).  For arm \(a\), multiply every normal-equation contribution by
\(\mathbf 1\{A_i=a\}/e_{i,a}(X_i)\).  Conditional on the observations preceding \(i\), its expectation is the corresponding full-data polynomial contribution, while its absolute magnitude is at most a fixed multiple of \(\varepsilon^{-1}\).  The density bounds make every population Gram matrix uniformly nonsingular.  Applying the scalar exponential-supermartingale bound
\[
\Pr\!\left(\left|\sum_{i=1}^m\xi_i\right|\geq x\right)
\leq 2\exp\!\left\{-\frac{x^2}{2(v+Bx/3)}\right\}
\tag{6}
\]
to a fixed net of the finite-dimensional basis, where \(|\xi_i|\leq B\) and the sum of the conditional variances is at most \(v\), shows that the empirical Gram matrix is singular or smaller than half its population lower bound with probability at most
\(C\exp(-cmb^d)\).  On that event use a fixed bounded fit.  Off that event, Taylor's formula gives integrated squared bias at most \(Cb^{2u}\), and direct conditional variance calculation gives \(C/(mb^d)\).  Projection of a fitted mean onto \([-M,M]\) cannot increase its pointwise error.  Consequently the mean estimator based on \(I_{j-2}\), with
\(b_{\mu,j}\asymp N_j^{-1/(2s+d)}\), satisfies
\[
\sup_{P\in\mathcal P(s,r)}
E_P\|\widehat\mu_{j,a}-\mu_a\|_{L_2(P_X)}^2
\leq C N_j^{-2\alpha_d}.
\tag{7}
\]
It is computed immediately after \(I_{j-2}\) and frozen before \(I_{j-1}\).

On \(I_{j-1}\), use the raw, untruncated response
\[
R_i:=\{Y_i-\widehat\mu_{j,a}(X_i)\}^2
\quad\text{when }A_i=a,
\tag{8}
\]
with the same recorded inverse-propensity weighting.  Conditional on the mean-training block,
\[
E_P[R_i\mid X_i=x,A_i=a]
=\sigma_a^2(x)+\{\widehat\mu_{j,a}(x)-\mu_a(x)\}^2.
\tag{9}
\]
No response clipping is performed, so (9) is exact.  Gaussianity, the variance range, and clipping only the fitted mean give a uniform conditional second moment for \(R_i\).  Let \(S_{j,a}\) be the degree-\(\lfloor r\rfloor\) local-polynomial smoother with
\(b_{\sigma,j}\asymp N_j^{-1/(2r+d)}\).  The same Gram-matrix calculation makes \(S_{j,a}\) uniformly bounded as an operator on \(L_2(P_X)\) on the good event.  Moreover,
\[
\begin{aligned}
E_P\|\{\widehat\mu_{j,a}-\mu_a\}^2\|_{L_2(P_X)}^2
&=E_P|\widehat\mu_{j,a}(X)-\mu_a(X)|^4\\
&\leq 4M^2E_P|\widehat\mu_{j,a}(X)-\mu_a(X)|^2.
\end{aligned}
\tag{10}
\]
Apply \(S_{j,a}\) to the raw responses (8), and only afterward project the fitted function onto \([v_{\min},v_{\max}]\); denote the result by \(\widehat v_{j,a}\).  Since projection cannot increase distance to \(\sigma_a^2\), equations (7)--(10), the \(r\)-Hölder bias bound, and the conditional variance calculation yield
\[
\sup_{P\in\mathcal P(s,r)}
E_P\|\widehat v_{j,a}-\sigma_a^2\|_{L_2(P_X)}^2
\leq C\left(N_j^{-2\beta_d}\vee N_j^{-2\alpha_d}\right).
\tag{11}
\]
These are expected integrated-risk bounds for one epoch.  No maximum over epochs and hence no union bound occurs.

At the start of \(I_j\), set
\[
\widehat z_{j,a}(x):=|\omega_a(x)|\sqrt{\widehat v_{j,a}(x)},
\qquad
e_{t,a}(x):=u_a^*(\widehat z_j(x)),\quad t\in I_j.
\tag{12}
\]
The rule is a law-independent measurable function of the current covariate and actually revealed past observations.  Equations (1) and (12) give normalization and \(e_{t,a}\geq\varepsilon\), so the resulting design \(D_n^{\circ}\) belongs to \(\mathfrak D_{n,\varepsilon}\).  Since both arguments of the square root are at least \(v_{\min}>0\),
\[
|\widehat z_{j,a}(x)-z_a(x)|^2
\leq \frac{W^2}{4v_{\min}}
|\widehat v_{j,a}(x)-\sigma_a^2(x)|^2.
\tag{13}
\]
The allocation-dependent part of \(V(e;P)\) is \(E_PF_{z(X)}\{e(X)\}\).  Integrating (4), and then using (11)--(13), proves, uniformly for \(t\in I_j\),
\[
E_P\{V(e_t;P)-V^*(P)\}
\leq C\left(N_j^{-2\beta_d}\vee N_j^{-2\alpha_d}\right).
\tag{14}
\]

Define the recorded-propensity score and estimator by
\[
\Psi_t:=
\sum_{a\in\mathcal A}\omega_a(X_t)\widehat\mu_{j,a}(X_t)
+\sum_{a\in\mathcal A}
\frac{\mathbf 1\{A_t=a\}}{e_{t,a}(X_t)}
\omega_a(X_t)\{Y_t-\widehat\mu_{j,a}(X_t)\},
\qquad
T_n^{\circ}:=\frac1n\sum_{t=1}^n\Psi_t,
\tag{15}
\]
for \(t\in I_j\); use any fixed bounded predictable initial fits and a fixed positive allocation on the initialization block.  Every fitted quantity in (15) is \(\mathcal G_{t-1}\)-measurable.  Conditional randomization, consistency, and positivity first give
\[
E[\Psi_t\mid\mathcal G_{t-1}^{+}]=h_P(X_t).
\]
Latent i.i.d. sampling and the fact that the construction excludes future baseline covariates then give
\[
E[\Psi_t\mid\mathcal G_{t-1}]=E_P h_P(X_t)=\theta(P).
\tag{16}
\]
Thus \(\Psi_t-\theta(P)\) is a square-integrable martingale difference for the smaller sequential-reveal filtration.  In particular, scores from distinct times are orthogonal even when the ambient \(H_0\) contains additional baseline covariates that the design ignores.

Put \(\delta_{j,a}:=\omega_a(\widehat\mu_{j,a}-\mu_a)\).  Conditional on \(\mathcal G_{t-1}^{+}\), the outcome residual has mean zero and variance \(\sigma_a^2(X_t)\).  Expanding the square in (15), then integrating over the fresh \(X_t\), gives the exact identity
\[
\begin{aligned}
E[(\Psi_t-\theta(P))^2\mid\mathcal G_{t-1}]
&=V(e_t;P)\\
&\quad+E_P\!\left[
\sum_{a\in\mathcal A}\frac{\delta_{j,a}(X)^2}{e_{t,a}(X)}
-\left\{\sum_{a\in\mathcal A}\delta_{j,a}(X)\right\}^2
\,\middle|\,\mathcal G_{t-1}\right].
\end{aligned}
\tag{17}
\]
The last square is nonnegative and every denominator is at least \(\varepsilon\).  Bounded weights and (7) therefore bound the remainder in (17) from above by \(CN_j^{-2\alpha_d}\).  Martingale orthogonality, followed by (14) and (17), now yields
\[
E[(T_n^{\circ}-\theta(P))^2]-\frac{V^*(P)}n
\leq \frac{C}{n^2}\sum_j |I_j|
\left(N_j^{-2\alpha_d}\vee N_j^{-2\beta_d}\right)
+\frac{C}{n^2}.
\tag{18}
\]
The last term covers the fixed initialization block.  Since \(|I_j|\leq N_j\), \(N_j=2^j\), and \(0<\alpha_d,\beta_d<1/2\), the geometric series in (18) is bounded by
\[
C\left(n^{1-2\alpha_d}\vee n^{1-2\beta_d}\right).
\]
After multiplication by \(n^{-2}\), this is the claimed
\(C(n^{-1-2\alpha_d}\vee n^{-1-2\beta_d})\) bound; the initialization term is smaller.  The bound is uniform over \(P\in\mathcal P(s,r)\), and (12) and (15) give one common data-driven pair rather than a law-indexed oracle pair.

Finally, if \(v_{\min}=v_{\max}\), every admissible conditional variance equals the known constant \(v_{\min}\).  In that boundary case the construction may set \(\widehat v_{j,a}\equiv v_{\min}\), so (14) is identically zero and the same proof gives the sharper bound \(Cn^{-1-2\alpha_d}\).  This refinement is not needed for the displayed claim but records the boundary behavior used by the nuisance-risk comparison.